\section*{Bab \ref{ch:g12:mountain-building} --- Pembentukan Pegunungan}

\begin{solution}{exo:g12:mountain-building:1}
(a) \omterm{def:g12:mountain-building:thrust}{sesar anjak}; (b) \omterm{def:g12:mountain-building:ophiolite}{ofiolit}; (c) \omterm{def:g12:mountain-building:fold}{sinklin}; (d) \omterm{def:g12:mountain-building:suture}{sutura}; (e) \omterm{def:g12:mountain-building:fold}{antiklin};
(f) \omterm{def:g12:mountain-building:nappe}{nappe}.
\end{solution}

\begin{solution}{exo:g12:mountain-building:2}
(1) Sayap-sayapnya miring menjauhi \omterm{def:g12:mountain-building:fold}{sumbu lipatan}: sebuah \omterm{def:g12:mountain-building:fold}{antiklin}. (2) Lapisan
yang tertua. (3) Tidak: erosi mengikuti lapisan yang keras dan yang lunak,
bukan bentuk \omterm{def:g12:mountain-building:fold}{lipatannya}; inti sebuah \omterm{def:g12:mountain-building:fold}{antiklin} dapat terkikis menjadi lembah.
\end{solution}

\begin{solution}{exo:g12:mountain-building:3}
(1) $e = 1 - \cos\ang{30} = 1 - 0.866 = 0.134$, sekitar \qty{13}{\percent}.
(2) $L_0 = 6.0/\cos\ang{30} \approx \qty{6.9}{km}$.
\end{solution}

\begin{solution}{exo:g12:mountain-building:4}
(1) \omterm{def:g11:ocean-ridges:normal-fault}{Sesar normal}; kerak itu teregang. (2) Ya: sesar itu adalah \omterm{def:g12:mountain-building:thrust}{sesar naik} yang
miring kurang dari \ang{45}.
\end{solution}

\begin{solution}{exo:g12:mountain-building:5}
(1) Sedimen laut, basalt bantal, \omterm{def:g11:igneous-rocks:intrusion}{retas} berlapis, \omterm{def:g11:ocean-ridges:gabbro}{gabro}, \omterm{def:g11:ocean-ridges:peridotite}{peridotit}. (2) Lava
membentuk bantal hanya bila mendingin di dalam air. (3) \omterm{def:g12:mountain-building:ophiolite}{Ofiolit} adalah
sepotong samudra yang dahulu terletak di antara dua benua; ketika kedua benua
itu bertumbukan, irisan-irisan dasarnya terdorong ke atas tepi-tepinya
alih-alih tenggelam, sehingga irisan-irisan itu dijumpai di tempat samudra itu
menutup, di sepanjang \omterm{def:g12:mountain-building:suture}{sutura}.
\end{solution}

\begin{solution}{exo:g12:mountain-building:6}
(1) $L_0 = 2.1 + 1.8 + 2.4 + 1.6 = \qty{7.9}{km}$;
$e = (7.9 - 6.0)/7.9 \approx 0.24$, \qty{24}{\percent}. (2) $\cos\theta =
6.0/7.9 \approx 0.759$, $\theta \approx \ang{41}$.
\end{solution}

\begin{solution}{exo:g12:mountain-building:7}
(1) $45\cos\ang{10} \approx \qty{44.3}{km}$. (2) $45\sin\ang{10} \approx
\qty{7.8}{km}$. (3) \omterm{def:g12:mountain-building:fold}{Lipatan} seperti itu memendek sebesar $1 - \cos\ang{12} \approx 0.0219$:
sebelum terlipat, sabuk itu pasti selebar $44.3/0.0219 \approx \qty{2000}{km}$.
\end{solution}

\begin{solution}{exo:g12:mountain-building:8}
(1) Superposisi: bila dibaca tanpa memperhitungkan sesar, batuan bagian atas
akan dianggap yang lebih muda. (2) $500 - 75 = 425$ juta tahun. (3) $\qty{40000}{m} /
\qty{800}{m} = 50$ kali.
\end{solution}

\begin{solution}{exo:g12:mountain-building:9}
(1) $(30.23 - 19.38) \times 111.2 = 10.85 \times 111.2 \approx
\qty{1207}{km}$. (2) $1207/14 \approx \qty{86}{km}$ per sejuta tahun,
\qty{86}{mm} atau sekitar \qty{8.6}{cm} setahun. (3) \qty{9}{m} seabad sama
dengan \qty{9}{cm} setahun: keduanya sesuai dengan selisih sekitar
\qty{5}{\percent}.
\end{solution}

\begin{solution}{exo:g12:mountain-building:10}
(1) $H = 35/(1 - 0.40) \approx \qty{58}{km}$. (2) $\Delta H \approx
\qty{23.3}{km}$; kerak itu berdiri lebih tinggi sebesar $23.3 \times 550/3380 \approx
\qty{3.8}{km}$, dan \omterm{def:g12:continental-crust:root}{akar} tambahannya sedalam $23.3 - 3.8 \approx \qty{19.5}{km}$.
\end{solution}

\begin{solution}{exo:g12:mountain-building:11}
(1) Gambar itu mengalikan $\tan\theta$ dengan 3: $\tan\theta = \tan\ang{40}/3
\approx 0.280$, $\theta \approx \ang{15.6}$. (2) Sebenarnya: $1 - \cos\ang{15.6}
\approx 0.037$, \qty{3.7}{\percent}; dibaca pada gambar: $1 - \cos\ang{40}
\approx 0.234$, \qty{23}{\percent}. (3) Skala vertikal yang diregangkan
mencuramkan kemiringan dan memperpanjang lapisan, sehingga \omterm{def:g12:mountain-building:shortening}{pemendekan} tampak
beberapa kali terlalu besar.
\end{solution}

\begin{solution}{exo:g12:mountain-building:12}
(1) $L_0 = 10/(1 - 0.317) \approx \qty{14.6}{km}$. (2) $1 - \cos\ang{57.5}
\approx 1 - 0.537 = 0.463$, \qty{46}{\percent}. (3) Pada \omterm{def:g12:mountain-building:fold}{lipatan} yang membulat
hanya sebagian kecil lapisan yang miring \ang{57.5}; di dekat engsel lapisan itu
hampir datar, dan bagian yang datar tidak memendek. Pada \omterm{def:g12:mountain-building:fold}{lipatan} bersayap lurus,
seluruh lapisan miring \ang{57.5}.
\end{solution}

\begin{solution}{exo:g12:mountain-building:13}
(1) $\sqrt{16.6^2 + 41.8^2} \approx \qty{45.0}{mm/yr}$. (2) $41.8 - 23.5 =
\qty{18.3}{mm}$ setahun, yaitu sekitar \qty{18}{km} \omterm{def:g12:mountain-building:shortening}{pemendekan} utara--selatan
per sejuta tahun di utara Lhasa. (3) Lhasa sendiri telah bergerak ke utara,
dan gempa-gempa bumi di sabuk itu tersebar jauh di utara Himalaya melintasi
Asia: kerak di sana masih terus patah.
\end{solution}

\begin{solution}{exo:g12:mountain-building:14}
(1) Kemiringannya hampir nol, sehingga \omterm{def:g12:mountain-building:shortening}{pemendekannya} sama dengan perpindahannya,
\qty{45}{km}. (2) $45/(1 - \cos\ang{12}) = 45/0.0219 \approx \qty{2060}{km}$.
(3) Konvergensi yang besar terutama diserap oleh \omterm{def:g12:mountain-building:thrust}{sesar anjak} dan \omterm{def:g12:mountain-building:nappe}{nappe} serta
oleh \omterm{def:g12:mountain-building:fold}{lipatan} yang ketat; \omterm{def:g12:mountain-building:fold}{lipatan} yang landai hanya menyerap sangat sedikit.
\end{solution}

\begin{solution}{exo:g12:mountain-building:15}
(1) $\qty{8848.86}{m} = \qty{8848860}{mm}$ dalam $\num{5e7}$ tahun: sekitar
\qty{0.18}{mm} setahun. (2) $40/0.18 \approx 220$: kenaikan itu lebih dari dua
ratus kali lebih lambat daripada gerak India. (3) Batu kapur itu diendapkan di
bawah permukaan laut, dan erosi telah menyingkirkan batuan yang dahulu berada di
atasnya: batu kapur itu telah terangkat lebih tinggi daripada ketinggiannya
sekarang.
\end{solution}

\begin{solution}{pb:g12:mountain-building:1}
\textbf{1.} $2\pi \times 6371/360 \approx \qty{111.2}{km}$.

\textbf{2.} Antara 50 juta tahun yang lalu (\ang{0.68}~LS) dan 45 (\ang{3.30}~LU):
$50 - 5 \times 0.68/3.98 \approx 49$ juta tahun yang lalu.

\textbf{3.} $(19.38 - 0.68) \times 111.2 = 18.70 \times 111.2 \approx
\qty{2079}{km}$.

\textbf{4.} $2079/16 \approx \qty{130}{km}$ per sejuta tahun: \qty{130}{mm}
setahun, sekitar \qty{13}{cm}.

\textbf{5.} 66 sampai 60: $7.28 \times 111.2/6 \approx 135$; 60 sampai 55: $4.86
\times 111.2/5 \approx 108$; 55 sampai 50: $6.56 \times 111.2/5 \approx 146$; 50
sampai 45: $3.98 \times 111.2/5 \approx 89$; 45 sampai 40: $2.29 \times 111.2/5
\approx 51$ (\unit{mm} setahun). Paling cepat dari 55 sampai 50 juta tahun yang
lalu.

\textbf{6.} $(20.00 - 5.59) \times 111.2/40 = 14.41 \times 111.2/40 \approx
\qty{40}{mm}$ setahun.

\textbf{7.} $40/130 \approx 0.31$: kurang dari sepertiga.

\textbf{8.} Dari langkah 55--50 (146) ke langkah 50--45 (89): sekitar 50 juta
tahun yang lalu.

\textbf{9.} \omterm{def:g12:continental-crust:continental-crust}{Kerak benua} terlalu ringan untuk tenggelam: ketika kerak itu
mencapai palung, \omterm{def:g11:subduction-zones:subduction}{subduksi} tidak lagi dapat menarik India dengan bebas, dan
benua yang lain menahannya.

\textbf{10.} $(20.00 + 0.68) \times 111.2 = 20.68 \times 111.2 \approx
\qty{2300}{km}$.

\textbf{11.} $(29.65 - 21.15) \times 111.2 = 8.50 \times 111.2 \approx
\qty{945}{km}$.

\textbf{12.} $2300 - 945 \approx \qty{1354}{km}$.

\textbf{13.} Kini $(29.65 - 20.00) \times 111.2 \approx \qty{1073}{km}$;
50 juta tahun yang lalu $(21.15 + 0.68) \times 111.2 = 21.83 \times 111.2
\approx \qty{2427}{km}$ (dan $2427 - 1073 = 1354$).

\textbf{14.} $e = 1354/2427 \approx 0.56$, \qty{56}{\percent}.

\textbf{15.} \omterm{def:g12:mountain-building:fold}{Lipatan}, dan \omterm{def:g12:mountain-building:thrust}{sesar naik}: \omterm{def:g12:mountain-building:thrust}{sesar anjak} yang membawa \omterm{def:g12:mountain-building:nappe}{nappe}.

\textbf{16.} $H = 35 \times 2427/1073 \approx \qty{79}{km}$.

\textbf{17.} $\Delta H = 79 - 35 = \qty{44}{km}$; kerak itu akan berdiri lebih
tinggi sebesar $44 \times 550/3380 \approx \qty{7.2}{km}$.

\textbf{18.} $\Delta H = 4.6 \times 3380/550 \approx \qty{28}{km}$, kerak setebal
$35 + 28 = \qty{63}{km}$, dekat dengan \qty{70}{km} yang
terukur.

\textbf{19.} Kerak itu tidak mempertahankan luas penampangnya: sebagian Asia
terjepit keluar ke arah timur, dan erosi menyingkirkan batuan dari puncak
pegunungan.

\textbf{20.} $2079 + 2300 \approx \qty{4380}{km}$: sejak akhir zaman Kapur,
India telah berpindah sekitar \textbf{\qty{4400}{km} ke utara}.
\end{solution}
